%=====================================================================
%  SECCION 3 -- MARCO NORMATIVO APLICABLE
%  Responsable: Fabian Moreno Ugarte
%
%  GENERADO desde secciones/03_marco_normativo.md, que sigue siendo la
%  FUENTE de esta seccion. No editar aqui: editar el .md y reconvertir.
%
%  Version comprimida (2026-09-14): el desarrollo completo esta en el
%  Anexo B, generado desde secciones/anexo_b_marco_normativo.md.
%  Los marcadores label sostienen las referencias cruzadas del resto
%  del informe. La numeracion 3.1-3.4 es vinculante.
%=====================================================================

\section{Marco normativo aplicable}
\label{sec:marco}

Esta sección identifica el régimen jurídico que rige el tratamiento de
imágenes de personas mediante videovigilancia y visión computacional en
un terminal aeroportuario, y traduce cada obligación en un requisito
verificable sobre el sistema propuesto. Las zonas en que la norma admite
más de una lectura razonable se presentan como riesgos a validar y no
como conclusiones: el análisis lo elabora un equipo de estudiantes de
ingeniería y no un despacho legal, y la calificación jurídica definitiva
corresponde al área legal de Lima Airport Partners y, en su caso, a la
Autoridad Nacional de Protección de Datos Personales. Aquí se recoge lo
esencial; el desarrollo completo, con el articulado citado literalmente,
figura en el Anexo B.

% ---- 3.1 ----
\subsection{Régimen normativo aplicable}
\label{sec:directiva}
\label{sec:dl1218}

El punto de partida es el artículo 2 de la Constitución, que en sus
incisos 6 y 7 protege la intimidad personal y familiar y la imagen
propia. Captar la imagen de una persona en un espacio abierto al público
no es por ello ilícito, pero constituye una injerencia en un derecho
fundamental y debe superar un examen de proporcionalidad.

La norma central es la Ley 29733, Ley de Protección de Datos Personales.
Su reglamento vigente es el Decreto Supremo 016-2024-JUS, publicado en
\emph{El Peruano} el 30 de noviembre de 2024, que derogó el D.S.
003-2013-JUS. Su entrada en vigor se habría producido el 30 de marzo de
2025, dato que se enuncia con reserva: es el único de los relativos a
este decreto que no se ha contrastado todavía contra la fuente oficial,
y es el que fija el régimen transitorio, de modo que debe confirmarse
antes de apoyarse en él. En cualquier caso, los procedimientos y
cláusulas que LAP tenga construidos bajo el reglamento anterior deben
revisarse. \verificar{únicamente la fecha de entrada en vigor, el 30 de
marzo de 2025, que es el dato que fija el régimen transitorio. La
expedición, la publicación y la disposición derogatoria sí se
contrastaron contra la separata de \emph{El Peruano} del 30 de noviembre
de 2024.}

Para videovigilancia rige además la Directiva 01-2020-JUS/DGTAIPD, que
sigue vigente: la disposición derogatoria del D.S. 016-2024-JUS deroga
únicamente el D.S. 003-2013-JUS. A ese cuerpo se superpone el régimen de
seguridad ciudadana del Decreto Legislativo 1218 y de la Ley 30120,
reglamentados por el Decreto Supremo 007-2020-IN. La denominación
oficial completa y la fecha de publicación del D.L. 1218 y de la Ley
30120 no se han contrastado todavía contra su texto publicado, a
diferencia del D.S. 007-2020-IN, por lo que ambas normas se citan por su
denominación usual y esos dos datos deben confirmarse antes de
reproducirse en un documento dirigido al cliente. El artículo 3 de ese
reglamento comprende los establecimientos comerciales abiertos al
público con un aforo de cincuenta personas o más, umbral que un terminal
aeroportuario supera con holgura: su aplicación al Aeropuerto Jorge
Chávez es el supuesto de partida y no una hipótesis remota.
\verificar{denominación oficial completa y fecha de publicación del D.L.
1218 y de la Ley 30120. El D.S. 007-2020-IN sí se contrastó contra el
texto publicado en \emph{El Peruano}.}

La zona gris principal es la calificación del propio terminal, un
espacio concesionado a un operador privado, abierto al público y con
funciones de interés público, que no encaja limpiamente ni como dominio
público ni como recinto privado. De esa calificación depende qué régimen
es exigible y con qué intensidad, y es el primer punto a elevar al área
legal del cliente.

% ---- 3.2 ----
\subsection{Datos tratados, legitimación y proporcionalidad}
\label{sec:proporcionalidad}

El sistema trata cuatro tipos de dato con regímenes distintos. El
fotograma contiene rostros identificables y es dato personal en el
sentido del artículo 2 numeral 4 de la Ley 29733. El punto con
coordenadas que producen los métodos basados en puntos
\cite{Song_2021_ICCV,Lin_2025_CVPR}, desligado del fotograma, no
identifica por sí solo a nadie. Asociado a una marca de tiempo y
encadenado con otros del mismo individuo, en cambio, reconstruye una
trayectoria, y esa diferencia depende de la política de persistencia y
no del algoritmo. El dato biométrico, por último, es sensible conforme
al numeral 5 solo cuando por sí mismo permite identificar al titular:
una plantilla facial lo es; un mapa de densidad o un conjunto de
coordenadas de cabezas, no. Que la salida del modelo no sea biométrica
no significa que no lo sea la entrada, de modo que almacenar o
transmitir fotogramas cambia el análisis aunque el modelo nunca ejecute
reconocimiento facial.

La presentación preliminar del proyecto habla de «personas
anonimizadas», y la documentación del prototipo califica su
identificador por persona de «no reversible». Un identificador
persistente que permite reagrupar las apariciones de una persona no es
anonimización, cuyo procedimiento es irreversible (numeral 12), sino
disociación, cuyo procedimiento es reversible (numeral 13). La
diferencia no es cosmética: el dato disociado sigue sujeto al régimen
general de la ley.

El consentimiento «previo, informado, expreso e inequívoco» del artículo
13.5 es inalcanzable en un terminal, porque no puede negarse sin
renunciar al servicio. La base habrá de buscarse entre los supuestos del
artículo 14, que la Directiva recoge en su numeral 6.3. El que la
práctica invoca con más frecuencia alude a los intereses legítimos del
titular de los datos, no a los del responsable como en el régimen
europeo, y de esa diferencia depende que ampare una finalidad operativa.
Qué base invoca LAP es información que solo el cliente puede aportar. La
finalidad, en cambio, no deja margen: los artículos 6 y 28 numeral 4
impiden extender el tratamiento a otra finalidad distinta de la
declarada. Medir la permanencia de personas en zonas comerciales, que el
propio proyecto llegó a contemplar como aplicación posible, no es una
variante de estimar aglomeraciones sino una finalidad distinta, con su
propia base de legitimación y su propio deber de información.

El artículo 7 exige que el tratamiento sea «adecuado, relevante y no
excesivo» para su finalidad. La estimación automática de densidad es
idónea, y si la finalidad es estimar aglomeraciones, el medio menos
invasivo que la alcanza produce conteos agregados sin identificar a
nadie. La ponderación en sentido estricto requiere, sin embargo,
evidencia sobre la magnitud real del problema que solo LAP puede
aportar, y mientras falte el examen descansa sobre una premisa no
acreditada. De este principio depende la decisión de mayor consecuencia
normativa: a igualdad de utilidad operativa, el mapa de densidad de
CSRNet \cite{Li_2018_CVPR} o la agregación por zona superan el examen de
necesidad mejor que la localización individual por puntos, cuya licitud
depende de que las coordenadas no se persistan ni se encadenen. La
elección del equipo no consta todavía por escrito.

% ---- 3.3 ----
\subsection{Privacidad desde el diseño y requisitos del sistema}
\label{sec:privacidad-diseno}

El D.S. 016-2024-JUS recoge en el Artículo IX de su Título Preliminar el
principio de responsabilidad proactiva: el responsable debe estar en
condiciones de acreditar que cumple, adoptando las medidas antes de
iniciar el tratamiento. De ahí se sigue la privacidad por diseño y por
defecto, que en este régimen es exigencia y no buena práctica, aunque la
norma no fija una lista mínima de medidas. El proyecto la traduce en
diez requisitos: procesamiento en el borde (M-01); no persistencia del
fotograma (M-02); salida agregada por zona, sin coordenadas individuales
persistidas (M-03); difuminado irreversible de rostros si algún
fotograma debe conservarse (M-04); plazo de retención con purga
automática (M-05); control de acceso por roles con auditoría (M-06);
cifrado en tránsito y en reposo (M-07); carteles informativos conforme
al Anexo A (M-08); linaje de datos documentado (M-09); y supervisión
humana, sin acciones automáticas sobre personas (M-10). Son requisitos y
no descripción del sistema actual: el prototipo respeta M-02, pero
\textbf{no cumple M-03}, como se expone en 3.4.

Sobre la retención, el artículo 8 de la Ley 29733 ordena conservar los
datos solo por el tiempo necesario, y la Directiva fija para las
imágenes un máximo de treinta a sesenta días. El artículo 17.2 del D.S.
007-2020-IN obliga, en cambio, a conservar las grabaciones un mínimo de
cuarenta y cinco días. Ambos regímenes solo son compatibles porque
recaen sobre finalidades distintas, y de ahí la consecuencia de diseño:
el sistema de conteo debe ser un subsistema separado del CCTV de
seguridad, con su propia finalidad, base y plazo, para no heredar una
obligación de conservación que no necesita.

Los artículos 18 a 25 reconocen los derechos de información, acceso,
rectificación, supresión, oposición, tutela e indemnización. El artículo
23, que protege frente a decisiones significativas basadas únicamente en
un tratamiento automatizado, es el fundamento directo de M-10: mientras
el sistema emita alertas agregadas que interpreta un operador, no se
activa. El deber de información se cumple mediante el cartel del numeral
6.11 de la Directiva, que junto con la proporcionalidad sostiene la
licitud del tratamiento; sus campos de identidad del responsable, base
de legitimación, canal de derechos y plazo de conservación solo puede
completarlos LAP.

Cualquier proveedor que trate datos por cuenta de LAP es encargado
sujeto al artículo 30. Si el procesamiento se ejecuta fuera del Perú,
hay flujo transfronterizo sujeto al artículo 15, y el Artículo VI del
Título Preliminar del reglamento extiende el régimen peruano a ciertos
tratamientos realizados en el exterior. Por eso la elección entre
inferencia \emph{on-premise} y en nube externa es, a juicio de este
informe, la de mayor impacto legal del proyecto; M-01 la anticipa en el
sentido más protector, pero la decisión no consta por escrito.

% ---- 3.4 ----
\subsection{Auditoría del prototipo: ¿el sistema trata datos biométricos?}
\label{sec:auditoria}
\label{sec:licencias}

La pregunta que decide el régimen aplicable es si en algún punto se
extrae de la imagen una característica que por sí sola permita
identificar a una persona. Bajo CSRNet \cite{Li_2018_CVPR} no se
construye ninguna plantilla y la salida es agregada por construcción.
Bajo P2PNet o P2R \cite{Song_2021_ICCV,Lin_2025_CVPR} tampoco hay
plantillas, pero la conformidad queda condicionada a que las coordenadas
no se persistan. Bajo el seguimiento multicámara con reidentificación
por apariencia la calificación cambia: tanto si el vector de apariencia
es dato sensible como si no lo es, el sistema reconstruye trayectorias
individuales y excede lo que la finalidad requiere conforme al artículo
7.

Contrastado con el prototipo que el equipo ha desarrollado, el resultado
es mixto. \textbf{La reidentificación por apariencia no se confirma.} La
fusión entre cámaras se resuelve por asociación espacio-temporal sobre
un plano común, y solo ante ambigüedad interviene un descriptor de
desempate (histograma de color del torso y proporción de la caja) que no
es biométrico y no se persiste. La decisión va en la dirección que exige
uno de los requisitos no funcionales que el proyecto se fijó, es decir,
de las condiciones que el sistema debe cumplir con independencia de su
funcionalidad: el que prohíbe derivar de la imagen rasgos que permitan
identificar a una persona, ni siquiera de forma indirecta (identificado
en el proyecto como RNF-01). Su texto literal y el documento en que
consta no figuran en la documentación disponible al cerrar esta versión,
de modo que se recoge según lo reportado por el equipo y esa conformidad
queda sujeta a contrastarlo con su fuente. \verificar{texto literal del
RNF-01 y documento en que consta. No figura en el repositorio al cerrar
esta versión, de modo que su formulación se recoge según lo reportado
por este frente y no según su fuente.} \textbf{Lo que se persiste, en
cambio, no es conforme.} El esquema de base de datos guarda, por cada
detección, un identificador de persona, la cámara, las coordenadas, el
instante y la zona: trayectorias individuales escritas en disco, que
incumplen M-03.

Del contraste resultan cuatro hallazgos no conformes. El primero es la
ausencia de la Evaluación de Impacto en Protección de Datos, exigible
antes de iniciar tratamientos de alto riesgo como la videovigilancia
masiva; es el más grave en lo formal y el más fácil de subsanar, porque
su contenido sustantivo está en la Sección 4. El segundo es el
incumplimiento de M-03, un defecto de implementación que el propio
equipo puede corregir agregando por zona antes de escribir o fijando una
retención muy corta sobre la tabla de posiciones. El tercero y el cuarto
son ausencias de evidencia que requieren acceso a material del terminal:
no hay validación de exactitud en el dominio de destino, y ninguno de
los trabajos revisados evalúa el sesgo demográfico de forma desagregada.
Quedan además pendientes la política de retención, la localización del
procesamiento y la trazabilidad de licencias de los conjuntos de datos.
